In isosceles triangle $OV_1V_2$, $\angle V_2V_1O=90^\circ-\dfrac{180^\circ}{n}$. The inscribed angle $\angle V_2V_1V_3$ equals $\dfrac{180^\circ}{n}$. The bisector condition therefore says $\dfrac{180}{n}=\dfrac{1}{2}(90-\dfrac{180}{n})$, yielding $\boxed{n=6}$. In a regular hexagon the ray to $V_3$ lies inside the specified angle and has these measures, verifying sufficiency.
\begin{asy}
import graph;import geometry;usepackage("amsmath");size(7cm);
path p;for(int k=0;k<6;++k)p=p--dir(60*k);draw(p--cycle);pair O=(0,0),A=dir(0),B=dir(60),C=dir(120);draw(A--O^^A--C,dashed);label("$V_1$",A,(1,0));label("$V_2$",B,(0,1));label("$V_3$",C,(0,1));label("$O$",O,(0,-1));
\end{asy}

